
%\newpage
\chapter{诗篇}
\section{卷一 1-41}
\subsection{第1篇：两条道路（義人与惡人道路的对比） 1:1-6}
\begin{table}[H]
\centering
\begin{tabular}{p{9.cm}|p{4.500cm}|p{2.7500cm}}\hline\hline
\multicolumn{1}{c|}{義人有福 1-3}&\multicolumn{1}{c|}{惡人滅亡 4-5}&\multicolumn{1}{c}{義人惡人的道路 6}\\\hline
不從惡人的計謀，不站罪人的道路，不坐褻慢人的座位，唯喜愛耶和華的律法，晝夜思想，這人便為有福！他要像一棵樹，栽在溪水旁，按時候結果子，葉子也不枯乾，凡他所做的，盡都順利。
&惡人並不是這樣，乃像糠秕，被風吹散。因此當審判的時候，惡人必站立不住，罪人在義人的會中也是如此。
&因為耶和華知道義人的道路，惡人的道路卻必滅亡。\\\hline
\end{tabular}
\end{table} 


\subsection{第2篇（\underline{基督的国度}）： 神膏立的君王必審判列國 2:1-12}
\begin{table}[H]
\centering
\begin{tabular}{p{3.75cm}|p{2.9cm}|p{4.7500cm}|p{4.7500cm}}\hline\hline
\multicolumn{1}{c|}{君民敵擋受膏者1-3}&\multicolumn{1}{c|}{主在錫安立君4-6}&\multicolumn{1}{c}{受膏者傳聖旨 7-9}&\multicolumn{1}{c}{悔改投靠基督 10-12}\\\hline
外邦為什麼爭鬧，萬民為什麼謀算虛妄的事？世上的君王一齊起來，臣宰一同商議，要抵擋耶和華並他的受膏者，說：「我們要掙開他們的捆綁，脫去他們的繩索！」
&那坐在天上的必發笑，主必嗤笑他們。那時，他要在怒中責備他們，在烈怒中驚嚇他們，說：「我已經立我的君在錫安我的聖山上了。」
&受膏者說：「我要傳聖旨。耶和華曾對我說：『你是我的兒子，我今日生你。你求我，我就將列國賜你為基業，將地極賜你為田產。你必用鐵杖打破他們，你必將他們如同窯匠的瓦器摔碎。』」
&現在你們君王應當醒悟，你們世上的審判官該受管教！當存畏懼侍奉耶和華，又當存戰兢而快樂。當以嘴親子，恐怕他發怒，你們便在道中滅亡，因為他的怒氣快要發作。凡投靠他的，都是有福的！\\\hline
\end{tabular}
\end{table} 


\subsection{第3篇（大衛逃避押沙龍）：在敌人攻击面前仍然信靠神的祷告 3:1-8}
\begin{table}[H]
\centering
\begin{tabular}{p{5.cm}|p{6.500cm}|p{5.2500cm}}\hline\hline
\multicolumn{1}{c|}{敌人的攻击 1-2}&\multicolumn{1}{c|}{神的看顾 3-6}&\multicolumn{1}{c}{大衛的祷告 7-8}\\\hline
大衛逃避他兒子押沙龍的時候作的詩。耶和華啊，我的敵人何其加增！有許多人起來攻擊我。有許多人議論我說：「他得不著神的幫助。」（細拉）
&但你耶和華是我四圍的盾牌，是我的榮耀，又是叫我抬起頭來的。我用我的聲音求告耶和華，他就從他的聖山上應允我。（細拉）我躺下睡覺，我醒著，耶和華都保佑我。雖有成萬的百姓來周圍攻擊我，我也不怕。
&耶和華啊，求你起來！我的神啊，求你救我！因為你打了我一切仇敵的腮骨，敲碎了惡人的牙齒。救恩屬乎耶和華，願你賜福給你的百姓。（細拉）\\\hline
\end{tabular}
\end{table} 


\subsection{第4篇（大衛）：上流虚妄人加羞辱时的祷告  4:1-8}
\begin{table}[H]
\centering
\begin{tabular}{p{3.cm}|p{8.500cm}|p{5.2500cm}}\hline\hline
\multicolumn{1}{c|}{在困苦中的祈祷1}&\multicolumn{1}{c|}{对上流虚妄人的劝诫  2-5}&\multicolumn{1}{c}{讲述神对自己的看顾  6-8}\\\hline
顯我為義的神啊，我呼籲的時候，求你應允我。我在困苦中，你曾使我寬廣；現在求你憐恤我，聽我的禱告。
&你們這上流人哪，你們將我的尊榮變為羞辱，要到幾時呢？你們喜愛虛妄，尋找虛假，要到幾時呢？（細拉）你們要知道，耶和華已經分別虔誠人歸他自己，我求告耶和華，他必聽我。你們應當畏懼，不可犯罪；在床上的時候要心裡思想，並要肅靜。（細拉）當獻上公義的祭，又當倚靠耶和華。
&有許多人說：「誰能指示我們什麼好處？」耶和華啊，求你仰起臉來，光照我們！你使我心裡快樂，勝過那豐收五穀、新酒的人。我必安然躺下睡覺，因為獨有你耶和華使我安然居住。\\\hline
\end{tabular}
\end{table} 


\subsection{第5篇（大衛）：刑罚以口犯罪的恶人，保护义人的祷告  5:1-12}
\subsubsection{诗人求神垂听祷告的呼求  1-3}
大衛的詩，交於伶長。用吹的樂器。
1 耶和華啊，求你留心聽我的言語，顧念我的心思。
2 我的王我的神啊，求你垂聽我呼求的聲音，因為我向你祈禱。
3 耶和華啊，早晨你必聽我的聲音；早晨我必向你陳明我的心意，並要警醒。

\subsubsection{态度 4-7}
\begin{table}[H]
\centering
\begin{tabular}{p{11.500cm}|p{5.2500cm}}\hline\hline
\multicolumn{1}{c|}{神对恶人的态度 4-6}&\multicolumn{1}{c}{自己对神的态度 7}\\\hline
因為你不是喜悅惡事的神，惡人不能與你同居。
5 狂傲人不能站在你眼前，凡作孽的都是你所恨惡的。
6 說謊言的，你必滅絕；好流人血弄詭詐的，都為耶和華所憎惡。
&至於我，我必憑你豐盛的慈愛進入你的居所，我必存敬畏你的心向你的聖殿下拜。\\\hline
\end{tabular}
\end{table} 

\subsubsection{讲述祷告的内容  8-12}
\begin{table}[H]
\centering
\begin{tabular}{p{2.9500cm}|p{4.00cm}|p{4.2500cm}|p{5.2500cm}}\hline\hline
\multicolumn{1}{c|}{求神引领走正路 8}&\multicolumn{1}{c|}{述说恶人口中的罪 9}&\multicolumn{1}{c|}{求神刑罚恶人 10}&\multicolumn{1}{c}{求神赐福保护义人 11-12}\\\hline
耶和華啊，求你因我的仇敵，憑你的公義引領我，使你的道路在我面前正直。
&因為他們的口中沒有誠實，他們的心裡滿有邪惡，他們的喉嚨是敞開的墳墓，他們用舌頭諂媚人。
&神啊，求你定他們的罪！願他們因自己的計謀跌倒，願你在他們許多的過犯中把他們逐出，因為他們背叛了你。
&凡投靠你的，願他們喜樂，時常歡呼，因為你護庇他們。又願那愛你名的人，都靠你歡欣。因為你必賜福於義人，耶和華啊，你必用恩惠如同盾牌四面護衛他。\\\hline
\end{tabular}
\end{table} 



\subsection{第6篇（大衛）：软弱病痛中的祷告  6:1-10}
\subsubsection{求神医治怜悯拯救的呼求 1-4}  
耶和華啊，求你不要在怒中責備我，也不要在烈怒中懲罰我！
2 耶和華啊，求你可憐我，因為我軟弱。耶和華啊，求你醫治我，因為我的骨頭發戰。
3 我心也大大地驚惶。耶和華啊，你要到幾時才救我呢？
4 耶和華啊，求你轉回，搭救我，因你的慈愛拯救我。  
\subsubsection{述说自己的愁苦及原因 5-7}  
5 因為在死地無人記念你，在陰間有誰稱謝你？
6 我因唉哼而困乏，我每夜流淚，把床榻漂起，把褥子濕透。我因憂愁眼睛乾癟，又因我一切的敵人眼睛昏花。
\subsubsection{相信神必使仇敌羞愧退後 8-10}
你們一切作孽的人，離開我吧！因為耶和華聽了我哀哭的聲音。
9 耶和華聽了我的懇求，耶和華必收納我的禱告。
10 我的一切仇敵都必羞愧，大大驚惶；他們必要退後，忽然羞愧。

\subsection{第7篇（大衛因便雅憫人古實的話）：求神向恶人施行公义审判的祷告  7:1-17}
\subsubsection{求神拯救的祈祷 1-10}
\begin{table}[H]
\centering
\begin{tabular}{p{33.em}|p{14.em}}\hline\hline
\multicolumn{1}{c|}{求神拯救脱离恶人1-6}&\multicolumn{1}{c}{求神在众人面前施行公义7-10}\\\hline
大衛指著便雅憫人古實的話，向耶和華唱的流離歌。
1 耶和華我的神啊，我投靠你，求你救我脫離一切追趕我的人，將我救拔出來，
2 恐怕他們像獅子撕裂我，甚至撕碎，無人搭救。
&\multirow{3}{14.em}{願眾民的會環繞你，願你從其上歸於高位。耶和華向眾民施行審判。耶和華啊，求你按我的公義和我心中的純正判斷我。願惡人的惡斷絕，願你堅立義人，因為公義的神察驗人的心腸肺腑。神是我的盾牌，他拯救心裡正直的人。}\\\cline{1-1}
3 耶和華我的神啊，我若行了這事，若有罪孽在我手裡，
4 我若以惡報那與我交好的人（連那無故與我為敵的，我也救了他），
5 就任憑仇敵追趕我，直到追上，將我的性命踏在地下，使我的榮耀歸於灰塵！（細拉）&\\\cline{1-1}
6 耶和華啊，求你在怒中起來，挺身而立，抵擋我敵人的暴怒；求你為我興起，你已經命定施行審判。&\\\hline
\end{tabular}
\end{table} 

\subsubsection{述说神向惡人的作为,稱謝神的公義  11-16}
\begin{table}[H]
\centering
\begin{tabular}{p{22.em}|p{19.5em}|p{6.5em}}\hline\hline
\multicolumn{1}{c|}{神待惡人的态度 11-13}&\multicolumn{1}{c|}{惡人的结局 14-16}&\multicolumn{1}{c}{对神的颂赞 17}\\\hline
11 神是公義的審判者，又是天天向惡人發怒的神。
12 若有人不回頭，他的刀必磨快，弓必上弦，預備妥當了。
13 他也預備了殺人的器械，他所射的是火箭。
&試看惡人因奸惡而劬勞，所懷的是毒害，所生的是虛假。
15 他掘了坑，又挖深了，竟掉在自己所挖的阱裡。
16 他的毒害必臨到他自己的頭上，他的強暴必落到他自己的腦袋上。
&我要照著耶和華的公義稱謝他，歌頌耶和華至高者的名。\\\hline
\end{tabular}
\end{table} 



\subsection{第8篇（\underline{基督的人性}大衛）：因神对人类将来的美好计划而赞美  8:1-9}
\begin{table}[H]
\centering
\begin{tabular}{p{10.em}|p{35em}|p{4.5em}}\hline\hline
\multicolumn{1}{c|}{赞美 1}&\multicolumn{1}{c|}{述说神对人类将来的美好计划  2-8}&\multicolumn{1}{c}{赞美 9}\\\hline
大衛的詩，交於伶長。用迦特樂器。耶和華我們的主啊，你的名在全地何其美，你將你的榮耀彰顯於天！
&你因敵人的緣故，從嬰孩和吃奶的口中建立了能力，使仇敵和報仇的閉口無言。我觀看你指頭所造的天，並你所陳設的月亮星宿，便說：「人算什麼，你竟顧念他？世人算什麼，你竟眷顧他？你叫他比天使(神)微小一點，並賜他榮耀尊貴為冠冕。你派他管理你手所造的，使萬物，就是一切的牛羊，田野的獸，空中的鳥，海裡的魚，凡經行海道的，都服在他的腳下。
&耶和華我們的主啊，你的名在全地何其美！\\\hline
\end{tabular}
\end{table} 


\subsection{第9篇（大衛）：称颂傳講神的審判的作為-憐憫,公義  1-20}
\subsubsection{要称颂传扬神的作为  1-10}
\begin{table}[H]
\centering
\begin{tabular}{p{14.em}|p{16em}|p{17.5em}}\hline\hline
\multicolumn{1}{c|}{四件我“要”1-3}&\multicolumn{1}{c|}{神已经完成的五个 3-6}&\multicolumn{1}{c}{四个神“要做”7-10}\\\hline
我\textcolor{red}{要}一心稱謝耶和華，我\textcolor{red}{要}傳揚你一切奇妙的作為。我\textcolor{red}{要}因你歡喜快樂，至高者啊，我\textcolor{red}{要}歌頌你的名！我的仇敵轉身退去的時候，他們一見你的面，\textcolor{red}{就}就跌倒滅亡。
&因你已經為我\textcolor{red}{申冤}，為我\textcolor{red}{辨屈}；你坐在寶座上，按公義審判。你曾\textcolor{red}{斥責}外邦，你曾\textcolor{red}{滅絕}惡人，你曾\textcolor{red}{塗抹}他們的名，直到永永遠遠。仇敵到了盡頭，他們被毀壞，直到永遠；你拆毀他們的城邑，連他們的名號都歸於無有。
&唯耶和華坐著\textcolor{red}{為王}，直到永遠，他已經為審判設擺他的寶座。他\textcolor{red}{要}按公義審判世界，按正直判斷萬民。耶和華又\textcolor{red}{要}給受欺壓的人做高臺，在患難的時候做高臺。耶和華啊，認識你名的人\textcolor{red}{要}倚靠你，因你\textcolor{red}{沒有}沒有離棄尋求你的人。\\\hline
\end{tabular}
\end{table} 

\subsubsection{诗人向神的四个祷告11-20}
\begin{table}[H]
\centering
\begin{tabular}{p{11.em}|p{28em}|p{8.5em}}\hline\hline
\multicolumn{1}{c|}{傳揚神的作为11-13a}&\multicolumn{1}{c|}{求神憐恤自己以及被神怜悯的结局13b-18}&\multicolumn{1}{c}{求神 7-10}\\\hline
應當歌頌居錫安的耶和華，將他所行的傳揚在眾民中！因為那追討流人血之罪的，他記念受屈的人，不忘記困苦人的哀求。耶和華啊，你是從死門把我提拔起來的，
&求你憐恤我，看那恨我的人所加給我的苦難，好叫我述說你一切的美德，我必在錫安城 （女子）的門因你的救恩歡樂。外邦人陷在自己所掘的坑中，他們的腳在自己暗設的網羅裡纏住了。耶和華已將自己顯明了，他已施行審判，惡人被自己手所做的纏住了（他叫惡人被自己手所做的累住了）。（細拉）惡人，就是忘記神的外邦人，都必歸到陰間。窮乏人必不永久被忘，困苦人的指望必不永遠落空。
&耶和華啊，求你起來，不容人得勝，願外邦人在你面前受審判！耶和華啊，求你使外邦人恐懼，願他們知道自己不過是人！（細拉）\\\hline
\end{tabular}
\end{table} 



\subsection{第10篇：述说恶人用诡计毒害人，求神解救伸冤的祷告  10:1-18}

\subsubsection{向神述說自己的困境和惡人的惡行 1-11}
\begin{table}[H]
\centering
\begin{tabular}{p{3.cm}|p{7.7500cm}|p{6.2500cm}}\hline\hline
\multicolumn{1}{c|}{自己的困境1-2}&\multicolumn{1}{c|}{驕傲以為沒有神3-7}&\multicolumn{1}{c}{殺害無辜,以為神忘記8-11}\\\hline
耶和華啊，你為什麼站在遠處？在患難的時候為什麼隱藏？
2 惡人在驕橫中把困苦人追得火急，願他們陷在自己所設的計謀裡。&因為惡人以心願自誇，貪財的背棄耶和華，並且輕慢他（他祝福貪財的，卻輕慢耶和華）。
4 惡人面帶驕傲，說：「耶和華必不追究。」他一切所想的，都以為沒有神。
5 凡他所做的，時常穩固，你的審判超過他的眼界。至於他一切的敵人，他都向他們噴氣。
6 他心裡說：「我必不動搖，世世代代不遭災難。」
7 他滿口是咒罵、詭詐、欺壓，舌底是毒害、奸惡。
&他在村莊埋伏等候，他在隱密處殺害無辜的人，他的眼睛窺探無倚無靠的人。
9 他埋伏在暗地，如獅子蹲在洞中。他埋伏，要擄去困苦人；他拉網，就把困苦人擄去。
10 他屈身蹲伏，無倚無靠的人就倒在他爪牙（強暴人）之下。
11 他心裡說：「神竟忘記了，他掩面永不觀看。」
\\\hline
\end{tabular}
\end{table} 

\subsubsection{求神解救受欺压的人，惩罚恶人的祷告 12-18}
\begin{table}[H]
\centering
\begin{tabular}{p{7.75cm}|p{4.00cm}|p{5.00cm}}\hline\hline
\multicolumn{1}{c|}{求神不任凭恶人12-14}&\multicolumn{1}{c|}{求神永做王，灭绝外邦15-16}&\multicolumn{1}{c}{求神垂听呼求，解救受欺人17-18}\\\hline
耶和華啊，求你起來！神啊，求你舉手，不要忘記困苦人！
13 惡人為何輕慢神，心裡說你必不追究？
14 其實你已經觀看，因為奸惡毒害你都看見了，為要以手施行報應。無倚無靠的人把自己交託你，你向來是幫助孤兒的。
&15 願你打斷惡人的膀臂，至於壞人，願你追究他的惡，直到淨盡。
16 耶和華永永遠遠為王，外邦人從他的地已經滅絕了。
&17 耶和華啊，謙卑人的心願你早已知道（聽見），你必預備他們的心，也必側耳聽他們的祈求，
18 為要給孤兒和受欺壓的人申冤，使強橫的人不再威嚇他們。\\\hline
\end{tabular}
\end{table} 



\subsection{第11篇（大衛）：?????上流虚妄人加羞辱时的祷告  11:1-7}
\begin{table}[H]
\centering
\begin{tabular}{p{4.75cm}|p{12.00cm}}\hline\hline
\multicolumn{1}{c|}{惡人 1-3}&\multicolumn{1}{c|}{試驗義人恨惡惡人 4-7}\\\hline
我是投靠耶和華，你們怎麼對我說：「你當像鳥，飛往你的山去！
看哪，惡人彎弓，把箭搭在弦上，要在暗中射那心裡正直的人。
3 根基若毀壞，義人還能做什麼呢？」
&4 耶和華在他的聖殿裡，耶和華的寶座在天上，他的慧眼察看世人。
5 耶和華試驗義人，唯有惡人和喜愛強暴的人，他心裡恨惡。
他要向惡人密布網羅，有烈火、硫磺、熱風做他們杯中的份。
7 因為耶和華是公義的，他喜愛公義，正直人必得見他的面。\\\hline
\end{tabular}
\end{table} 

\begin{itemize}
\itemsep-.45em  
\item 在困苦中的祈祷  1
\item 对上流虚妄人的劝诫  2-5
\item 讲述神对自己的看顾  6-8
\end{itemize}

\subsection{第12篇（大衛）：求神剪除虚假和谎言  12:1-8}
\begin{table}[H]
\centering
\begin{tabular}{p{4.75cm}|p{7.00cm}|p{4.00cm}}\hline\hline
\multicolumn{1}{c|}{求神剪除虚假和谎言祈祷   1-3}&\multicolumn{1}{c|}{恶人和神的言语  4-6}&\multicolumn{1}{c|}{求神保守脱离恶人  7-8}\\\hline
耶和華啊，求你幫助，因虔誠人斷絕了，世人中間的忠信人沒有了。
2 人人向鄰舍說謊，他們說話是嘴唇油滑，心口不一。
3 凡油滑的嘴唇和誇大的舌頭，耶和華必要剪除。
&他們曾說：「我們必能以舌頭得勝，我們的嘴唇是我們自己的，誰能做我們的主呢？」
5 耶和華說：「因為困苦人的冤屈和貧窮人的嘆息，我現在要起來，把他安置在他所切慕的穩妥之地。」耶和華的言語是純淨的言語，如同銀子在泥爐中煉過七次。
&耶和華啊，你必保護他們，你必保佑他們永遠脫離這世代的人。
8 下流人在世人中升高，就有惡人到處遊行。\\\hline
\end{tabular}
\end{table} 


\subsection{第13篇（大衛）：在憂苦中祈求神看顧 13:1-6}
\begin{table}[H]
\centering
\begin{tabular}{p{12.cm}|p{5.00cm}}\hline\hline
\multicolumn{1}{c|}{在憂苦中祈求神 1-4}&\multicolumn{1}{c}{要赞美神的救恩 5-6}\\\hline
耶和華啊，你忘記我要到幾時呢？要到永遠嗎？你掩面不顧我要到幾時呢？
2 我心裡籌算，終日愁苦，要到幾時呢？我的仇敵升高壓制我，要到幾時呢？
3 耶和華我的神啊，求你看顧我，應允我，使我眼目光明，免得我沉睡至死，
4 免得我的仇敵說「我勝了他」，免得我的敵人在我搖動的時候喜樂。
&但我倚靠你的慈愛，我的心因你的救恩快樂。
6 我要向耶和華歌唱，因他用厚恩待我。\\\hline
\end{tabular}
\end{table} 


\subsection{第14篇（大衛）：神救護子民脫離愚頑世人  14:1-7}
\begin{table}[H]
\centering
\begin{tabular}{p{3.cm}|p{7.00cm}|p{7.00cm}}\hline\hline
\multicolumn{1}{c|}{愚頑世人1}&\multicolumn{1}{c}{神对恶人的态度2-4}&\multicolumn{1}{c}{神对義人的态度5-7}\\\hline
愚頑人心裡說：「沒有神！」他們都是邪惡，行了可憎惡的事，沒有一個人行善。
&耶和華從天上垂看世人，要看有明白的沒有，有尋求神的沒有。
3 他們都偏離正路，一同變為汙穢；並沒有行善的，連一個也沒有。
4 作孽的都沒有知識嗎？他們吞吃我的百姓，如同吃飯一樣，並不求告耶和華。
&他們在那裡大大地害怕，因為神在義人的族類中。
6 你們叫困苦人的謀算變為羞辱，然而耶和華是他的避難所。
7 但願以色列的救恩從錫安而出！耶和華救回他被擄的子民，那時雅各要快樂，以色列要歡喜。\\\hline
\end{tabular}
\end{table} 


\subsection{第15篇（大衛）： 得居聖山者之品行 15:1-5}
耶和華啊，誰能寄居你的帳幕？誰能住在你的聖山？
2 就是行為正直，做事公義，心裡說實話的人。
3 他不以舌頭讒謗人，不惡待朋友，也不隨夥毀謗鄰里。
4 他眼中藐視匪類，卻尊重那敬畏耶和華的人。他發了誓，雖然自己吃虧，也不更改。
5 他不放債取利，不受賄賂以害無辜。行這些事的人，必永不動搖。

\subsection{第16篇（\underline{基督的复活}大衛）：耶和華是我的產業,福樂并满足  16:1-11}
\subsubsection{我一心投靠耶和华 1-6}
\begin{table}[H]
\centering
\begin{tabular}{p{12.cm}|p{5.cm}}\hline\hline
\multicolumn{1}{c|}{我对待耶和华}&\multicolumn{1}{c}{耶和华对待我}\\\hline
神啊，求你保佑我，因為我投靠你我的心哪，你曾對耶和華說：「你是我的主，我的好處不在你以外。」論到世上的聖民，他們又美又善，是我最喜悅的。以別神代替耶和華的(送禮物給別神的)，他們的愁苦必加增。他們所澆奠的血我不獻上，我嘴唇也不提別神的名號。
&耶和華是我的產業，是我杯中的份，我所得的你為我持守。用繩量給我的地界坐落在佳美之處，我的產業實在美好！\\\hline
\end{tabular}
\end{table} 
\subsubsection{我靠耶和华欢喜并满足 7-11}
\begin{table}[H]
\centering
\begin{tabular}{p{9.5cm}|p{7.24cm}}\hline\hline
\multicolumn{1}{c|}{我对待耶和华}&\multicolumn{1}{c}{耶和华对待我}\\\hline
我必稱頌那指教我的耶和華，我的心腸在夜間也警戒我。我將耶和華常擺在我面前，因他在我右邊，我便不致搖動。因此我的心歡喜，我的靈(榮耀)快樂，我的肉身也要安然居住。
&因為你必不將我的靈魂撇在陰間，也不叫你的聖者見朽壞。你必將生命的道路指示我，在你面前有滿足的喜樂，在你右手中有永遠的福樂。\\\hline
\end{tabular}
\end{table} 

\subsection{第17（大衛）篇： 求神得脫惡敵  17:1-15}
\subsubsection{求神聽聞公義 1-5}
\begin{table}[H]
\centering
\begin{tabular}{p{7.cm}|p{11.cm}}\hline\hline
\multicolumn{1}{c|}{三聽二願}&\multicolumn{1}{c}{被神鑒察无過}\\\hline
耶和華啊，求你聽聞公義，側耳聽我的呼籲，求你留心聽我這不出於詭詐嘴唇的祈禱。願我的判語從你面前發出，願你的眼睛觀看公正。
&你已經試驗我的心，你在夜間鑒察我，你熬煉我，卻找不著什麼。我立志叫我口中沒有過失。論到人的行為，我藉著你嘴唇的言語，自己謹守，不行強暴人的道路。我的腳踏定了你的路徑，我的兩腳未曾滑跌。\\\hline
\end{tabular}
\end{table} 
 
\subsubsection{我靠耶和华欢喜并满足 6-15}
\begin{table}[H]
\centering
\begin{tabular}{p{5.9cm}|p{12.cm}}\hline\hline
\multicolumn{1}{c|}{求告神必應允}&\multicolumn{1}{c}{惡人的结局}\\\hline
\multirow{2}{6cm}{神啊，我曾求告你，因為你必應允我，求你向我側耳，聽我的言語。求你顯出你奇妙的慈愛來，你是那用右手拯救投靠你的，脫離起來攻擊他們的人。求你保護我，如同保護眼中的瞳人，將我隱藏在你翅膀的蔭下，使我脫離那欺壓我的惡人，就是圍困我，要害我命的仇敵}
&他們的心被脂油包裹，他們用口說驕傲的話。他們圍困了我們的腳步，他們瞪著眼，要把我們推倒在地。他像獅子急要抓食，又像少壯獅子蹲伏在暗處。耶和華啊，求你起來，前去迎敵，將他打倒，用你的刀救護我命脫離惡人。耶和華啊，求你用手救我脫離世人，脫離那只在今生有福分的世人。你把你的財寶充滿他們的肚腹，他們因有兒女就心滿意足，將其餘的財物留給他們的嬰孩。\\\cline{2-2}
&至於我，我必在義中見你的面；我醒了的時候得見（看）你的形象，就心滿意足了。\\\hline
\end{tabular}
\end{table} 



\subsection{第18篇（大衛脱离扫罗）： 因勝仇敵做王稱頌主恩 18:1-50}
\subsubsection{在急難中求告耶和華,神救我脱离仇敌 1-18}
\begin{table}[H]
\centering
\begin{tabular}{p{7.25cm}|p{7.5cm}|p{2.5cm}}\hline\hline
\multicolumn{1}{c|}{求告1-6}&\multicolumn{1}{c|}{神的威严7-15}&\multicolumn{1}{c}{神救我16-18}\\\hline
耶和華的僕人大衛的詩，交於伶長。當耶和華救他脫離一切仇敵和掃羅之手的日子，他向耶和華念這詩的話，說：
耶和華我的力量啊，我愛你！ 耶和華是我的巖石，我的山寨，我的救主，我的神，我的磐石，我所投靠的。他是我的盾牌，是拯救我的角，是我的高臺。
我要求告當讚美的耶和華，這樣我必從仇敵手中被救出來。曾有死亡的繩索纏繞我，匪類的急流使我驚懼，陰間的繩索纏繞我，死亡的網羅臨到我。我在急難中求告耶和華，向我的神呼求。他從殿中聽了我的聲音，我在他面前的呼求入了他的耳中。
&7 那時因他發怒，地就搖撼戰抖，山的根基也震動搖撼。
8 從他鼻孔冒煙上騰，從他口中發火焚燒，連炭也著了。
9 他又使天下垂，親自降臨，有黑雲在他腳下。
10 他坐著基路伯飛行，他藉著風的翅膀快飛。
11 他以黑暗為藏身之處，以水的黑暗、天空的厚雲為他四圍的行宮。
12 因他面前的光輝，他的厚雲行過，便有冰雹火炭。
13 耶和華也在天上打雷，至高者發出聲音，便有冰雹火炭。
14 他射出箭來，使仇敵四散；多多發出閃電，使他們擾亂。
15 耶和華啊，你的斥責一發，你鼻孔的氣一出，海底就出現，大地的根基也顯露。
&他從高天伸手抓住我，把我從大水中拉上來。17 他救我脫離我的勁敵和那些恨我的人，因為他們比我強盛。18 我遭遇災難的日子，他們來攻擊我，但耶和華是我的倚靠。\\\hline
\end{tabular}
\end{table} 

\subsubsection{神領我到寬闊之處,遵守神典章明白神的心意，心中被燈照明 19-28}
他又領我到寬闊之處；他救拔我，因他喜悅我。20 耶和華按著我的公義報答我，按著我手中的清潔賞賜我。21 因為我遵守了耶和華的道，未曾作惡離開我的神。22 他的一切典章常在我面前，他的律例我也未曾丟棄。23 我在他面前做了完全人，我也保守自己遠離我的罪孽。24 所以耶和華按我的公義，按我在他眼前手中的清潔償還我。
25 慈愛的人，你以慈愛待他；完全的人，你以完全待他。
26 清潔的人，你以清潔待他；乖僻的人，你以彎曲待他。
27 困苦的百姓，你必拯救；高傲的眼目，你必使他降卑。
你必點著我的燈，耶和華我的神必照明我的黑暗。


\subsubsection{藉著神衝入敵軍,神教導我能以爭戰,外邦人要投降 29-45}
\begin{table}[H]
\centering
\begin{tabular}{p{7.25cm}|p{6.cm}|p{4.05cm}}\hline\hline
\multicolumn{1}{c|}{衝入敵軍29-36}&\multicolumn{1}{c|}{追趕我的仇敵37-41}&\multicolumn{1}{c}{得胜42-45}\\\hline
29 我藉著你衝入敵軍，藉著我的神跳過牆垣。
30 至於神，他的道是完全的，耶和華的話是煉淨的，凡投靠他的，他便做他們的盾牌。
31 除了耶和華，誰是神呢？除了我們的神，誰是磐石呢？
32 唯有那以力量束我的腰，使我行為完全的，他是神。
33 他使我的腳快如母鹿的蹄，又使我在高處安穩。
34 他教導我的手能以爭戰，甚至我的膀臂能開銅弓。
35 你把你的救恩給我做盾牌，你的右手扶持我，你的溫和使我為大。
36 你使我腳下的地步寬闊，我的腳未曾滑跌。
&我要追趕我的仇敵，並要追上他們，不將他們滅絕，我總不歸回。
我要打傷他們，使他們不能起來，他們必倒在我的腳下。
因為你曾以力量束我的腰，使我能爭戰，你也使那起來攻擊我的都服在我以下。
你又使我的仇敵在我面前轉背逃跑，叫我能以剪除那恨我的人。
他們呼求，卻無人拯救；就是呼求耶和華，他也不應允。
&我搗碎他們，如同風前的灰塵；倒出他們，如同街上的泥土。
你救我脫離百姓的爭競，立我做列國的元首，我素不認識的民必侍奉我。
他們一聽見我的名聲就必順從我，外邦人要投降我。
外邦人要衰殘，戰戰兢兢地出他們的營寨。\\\hline
\end{tabular}
\end{table} 

\subsubsection{赞美神  46-50}
耶和華是活神！願我的磐石被人稱頌，願救我的神被人尊崇！
47 這位神就是那為我申冤，使眾民服在我以下的。
48 你救我脫離仇敵，又把我舉起，高過那些起來攻擊我的，你救我脫離強暴的人。
耶和華啊，因此我要在外邦中稱謝你，歌頌你的名。耶和華賜極大的救恩給他所立的王，施慈愛給他的受膏者，就是給大衛和他的後裔，直到永遠。



\subsection{第19篇（大衛）： 神荣耀的创造，全备的律法，赦罪之恩 19:1-14}
\subsubsection{神的创造彰显神的荣耀 1-6}
大衛的詩，交於伶長。
1 諸天述說神的榮耀，穹蒼傳揚他的手段。
2 這日到那日發出言語，這夜到那夜傳出知識。
3 無言無語，也無聲音可聽。
4 他的量帶通遍天下，他的言語傳到地極。神在其間為太陽安設帳幕，
5 太陽如同新郎出洞房，又如勇士歡然奔路。
6 它從天這邊出來，繞到天那邊，沒有一物被隱藏不得他的熱氣。
\subsubsection{神的律法使人完全 7-11}
\begin{table}[H]
\centering
\begin{tabular}{p{6.65cm}|p{10.25cm}}\hline\hline
\multicolumn{1}{c|}{神律法的功用}&\multicolumn{1}{c}{神律法的宝贵}\\\hline
耶和華的律法全備，能甦醒人心。耶和華的法度確定，能使愚人有智慧。
8 耶和華的訓詞正直，能快活人的心。耶和華的命令清潔，能明亮人的眼目。
&耶和華的道理潔淨，存到永遠。耶和華的典章真實，全然公義。
10 都比金子可羨慕，且比極多的精金可羨慕；比蜜甘甜，且比蜂房下滴的蜜甘甜。
11 況且你的僕人因此受警戒，守著這些便有大賞。\\\hline
\end{tabular}
\end{table} 

\subsubsection{求神救我脱离罪恶 12-14}
誰能知道自己的錯失呢？願你赦免我隱而未現的過錯。
13 求你攔阻僕人，不犯任意妄為的罪，不容這罪轄制我，我便完全，免犯大罪。
14 耶和華我的磐石，我的救贖主啊，願我口中的言語、心裡的意念在你面前蒙悅納。


\subsection{第20篇（大衛）： 祈耶和華援其受膏者使之得勝 20:1-9}
\begin{table}[H]
\centering
\begin{tabular}{p{6.65cm}|p{10.25cm}}\hline\hline
\multicolumn{1}{c|}{向神的心愿 1-5}&\multicolumn{1}{c}{耶和華援其受膏者 5-9}\\\hline
大衛的詩，交於伶長。
1 願耶和華在你遭難的日子應允你，願名為雅各神的高舉你。
2 願他從聖所救助你，從錫安堅固你；
3 記念你的一切供獻，悅納你的燔祭；（細拉）
4 將你心所願的賜給你，成就你的一切籌算。
5 我們要因你的救恩誇勝，要奉我們神的名豎立旌旗。願耶和華成就你一切所求的！
&現在我知道耶和華救護他的受膏者，必從他的聖天上應允他，用右手的能力救護他。
7 有人靠車，有人靠馬，但我們要提到耶和華我們神的名。
8 他們都屈身仆倒，我們卻起來，立得正直。
9 求耶和華施行拯救，我們呼求的時候，願王應允我們！\\\hline
\end{tabular}
\end{table} 


\subsection{第21篇（大衛）： 歌頌耶和華对王的救恩，和对仇敌的报应 21:1-13}
\begin{table}[H]
\centering
\begin{tabular}{p{8.25cm}|p{6.75cm}|p{2.cm}}\hline\hline
\multicolumn{1}{c|}{神对王的赐福 1-7}&\multicolumn{1}{c}{神对仇敌的报应 8-12}&\multicolumn{1}{c}{歌頌13}\\\hline
大衛的詩，交於伶長。
1 耶和華啊，王必因你的能力歡喜，因你的救恩他的快樂何其大！
2 他心裡所願的，你已經賜給他；他嘴唇所求的，你未嘗不應允。（細拉）
3 你以美福迎接他，把精金的冠冕戴在他頭上。
4 他向你求壽，你便賜給他，就是日子長久，直到永遠。
5 他因你的救恩大有榮耀，你又將尊榮威嚴加在他身上。
6 你使他有洪福，直到永遠，又使他在你面前歡喜快樂。
7 王倚靠耶和華，因至高者的慈愛，必不搖動。
&你的手要搜出你的一切仇敵，你的右手要搜出那些恨你的人。
9 你發怒的時候，要使他們如在炎熱的火爐中。耶和華要在他的震怒中吞滅他們，那火要把他們燒盡了。
10 你必從世上滅絕他們的子孫（果子），從人間滅絕他們的後裔。
11 因為他們有意加害於你，他們想出計謀，卻不能做成。
12 你必使他們轉背逃跑，向他們的臉搭箭在弦。
&耶和華啊，願你因自己的能力顯為至高！這樣，我們就唱詩，歌頌你的大能。\\\hline
\end{tabular}
\end{table} 




\subsection{第22篇  （\underline{基督的受苦得荣} 大衛調用朝鹿） ：  22:1-31}
个人哀歌, 同时也是弥赛亚诗篇,记述弥赛亚的受难与复活。 彼前1:10-11, “论到这救恩，那豫先说你们要得恩典的众先知，早已详细的寻求考察。就是考察在他们心里基督的灵，豫先证明基督受苦难，后来得荣耀，是指着什么时候，并怎样的时候。”
\begin{itemize}
\itemsep-.45em  
\item 我: 22:1-2;   你: 22:3-5;         
\item 我: 22:6-8;   你: 22:9-11;
\item 我: 22:12-18; 你: 22:19-21;
\item 我: 22:22-25; 你:26:31
\end{itemize}

\subsubsection{在困境中呼求神帮助 - \underline{基督受苦} 1-21}
\begin{table}[H]
\centering
\begin{tabular}{p{8.25cm}|p{4.25cm}|p{5.cm}}\hline\hline
\multicolumn{1}{c|}{自己被人羞辱藐视求神帮助 1-10}&\multicolumn{1}{c}{仇敌的围困受惊惧 11-15}&\multicolumn{1}{c}{生命受到威胁，求神拯救16-21}\\\hline
大衛的詩，交於伶長。調用朝鹿。
1 我的神！我的神！為什麼離棄我？為什麼遠離不救我，不聽我唉哼的言語？
2 我的神啊，我白日呼求，你不應允；夜間呼求，並不住聲。
3 但你是聖潔的，是用以色列的讚美為寶座(居所)的。
4 我們的祖宗倚靠你，他們倚靠你，你便解救他們。
5 他們哀求你便蒙解救，他們倚靠你就不羞愧。
6 但我是蟲，不是人，被眾人羞辱，被百姓藐視。
7 凡看見我的都嗤笑我，他們撇嘴搖頭，說：
8 「他把自己交託耶和華，耶和華可以救他吧！耶和華既喜悅他，可以搭救他吧！」
9 但你是叫我出母腹的，我在母懷裡，你就使我有倚靠的心。
10 我自出母胎就被交在你手裡，從我母親生我，你就是我的神。
&求你不要遠離我，因為急難臨近了，沒有人幫助我。
12 有許多公牛圍繞我，巴珊大力的公牛四面困住我。
13 牠們向我張口，好像抓撕吼叫的獅子。
14 我如水被倒出來，我的骨頭都脫了節，我心在我裡面如蠟熔化。
15 我的精力枯乾，如同瓦片，我的舌頭貼在我牙床上。你將我安置在死地的塵土中。
&犬類圍著我，惡黨環繞我，他們扎了我的手、我的腳。
17 我的骨頭我都能數過，他們瞪著眼看我。
18 他們分我的外衣，為我的裡衣拈鬮。
19 耶和華啊，求你不要遠離我！我的救主啊，求你快來幫助我！
20 求你救我的靈魂脫離刀劍，救我的生命(獨一者)脫離犬類！
21 救我脫離獅子的口！你已經應允我，使我脫離野牛的角。\\\hline
\end{tabular}
\end{table} 


\subsubsection{讚美主的恩佑 - \underline{基督得荣}  22-31}
\begin{itemize}
\itemsep-.45em  
\item 我要將你的名傳於我的弟兄，在會中我要讚美你 22-24
\begin{table}[H]
\centering
\begin{tabular}{p{10.65cm}|p{7.25cm}}\hline\hline
\multicolumn{1}{c|}{赞美的人群}&\multicolumn{1}{c}{赞美的原因}\\\hline
我要將你的名傳於我的弟兄，在會中我要讚美你。
23 你們敬畏耶和華的人，要讚美他！雅各的後裔，都要榮耀他！以色列的後裔，都要懼怕他！
&因為他沒有藐視、憎惡受苦的人，也沒有向他掩面，那受苦之人呼籲的時候，他就垂聽。\\\hline
\end{tabular}
\end{table} 


\item 我在大會中讚美你的話是從你而來的，我要在敬畏耶和華的人面前還我的願 25-31
\begin{table}[H]
\centering
\begin{tabular}{p{13cm}|p{4.5cm}}\hline\hline
\multicolumn{1}{c|}{讚美你的話25-28}&\multicolumn{1}{c}{還我的願30-31}\\\hline
我在大會中讚美你的話是從你而來的，我要在敬畏耶和華的人面前還我的願。
26 謙卑的人必吃得飽足，尋求耶和華的人必讚美他。願你們的心永遠活著！
27 地的四極都要想念耶和華，並且歸順他，列國的萬族都要在你面前敬拜。
28 因為國權是耶和華的，他是管理萬國的。
29 地上一切豐肥的人必吃喝而敬拜，凡下到塵土中不能存活自己性命的人，都要在他面前下拜。
&他必有後裔侍奉他，主所行的事必傳於後代。
31 他們必來把他的公義傳給將要生的民，言明這事是他所行的。\\\hline
\end{tabular}
\end{table} 

\end{itemize}


\subsection{第23篇（\underline{基督是牧者}大衛）：  耶和華是我的牧者 23:1-6}
耶和華是我的牧者，我必不致缺乏。2 他使我躺臥在青草地上，領我在可安歇的水邊。3 他使我的靈魂甦醒，為自己的名引導我走義路。4 我雖然行過死蔭的幽谷，也不怕遭害，因為你與我同在，你的杖、你的竿都安慰我。5 在我敵人面前，你為我擺設筵席，你用油膏了我的頭，使我的福杯滿溢。6 我一生一世必有恩惠、慈愛隨著我，我且要住在耶和華的殿中，直到永遠。

\subsection{第24篇（\underline{基督是君王}大衛）： 神在拣选的城中作荣耀的王 24:1-10}
\subsubsection{蒙耶和華賜福并成義的族类 1-6}
\begin{table}[H]
\centering
\begin{tabular}{p{8cm}|p{8.5cm}}\hline\hline
\multicolumn{1}{c|}{神大能的创造1-3}&\multicolumn{1}{c}{登耶和华圣山的人4-6}\\\hline
地和其中所充滿的，世界和住在其間的，都屬耶和華。
2 他把地建立在海上，安定在大水之上。
3 誰能登耶和華的山？誰能站在他的聖所？
&就是手潔心清，不向虛妄，起誓不懷詭詐的人。
5 他必蒙耶和華賜福，又蒙救他的神使他成義。
6 這是尋求耶和華的族類，是尋求你面的雅各。（細拉）\\\hline
\end{tabular}
\end{table} 

\subsubsection{眾城門抬起頭來，迎接榮耀的王進來 7-10}
眾城門哪，你們要抬起頭來！永久的門戶，你們要被舉起！那榮耀的王將要進來！
8 榮耀的王是誰呢？就是有力有能的耶和華，在戰場上有能的耶和華。
9 眾城門哪，你們要抬起頭來！永久的門戶，你們要把頭抬起！那榮耀的王將要進來！
10 榮耀的王是誰呢？萬軍之耶和華，他是榮耀的王。（細拉）


\subsection{第25篇（大衛）： 诗人在患难中呼求神的怜悯和引导，并述说神的恩典 25:1-22}
\subsubsection{诗人向神的祈求 1-7}
\begin{table}[H]
\centering
\begin{tabular}{p{5.75cm}|p{5.00cm}|p{6.00cm}}\hline\hline
\multicolumn{1}{c|}{两个“不羞愧”一个“必”1-3}&\multicolumn{1}{c|}{两个“教訓引導”4-5}&\multicolumn{1}{c}{两“念”一“不念”6-7}\\\hline
耶和華啊，我的心仰望你。
2 我的神啊，我素來倚靠你，求你不要叫我\textcolor{red}{羞愧}，不要叫我的仇敵向我誇勝。
3 凡等候你的必不\textcolor{red}{羞愧}，唯有那無故行奸詐的必要\textcolor{red}{羞愧}。
&耶和華啊，求你將你的道指示我，將你的路\textcolor{red}{教訓}我。5 求你以你的真理引導我，\textcolor{red}{教訓}我，因為你是救我的神，我終日等候你。
&耶和華啊，求你\textcolor{red}{記念}你的憐憫和慈愛，因為這是亙古以來所常有的。
7 求你\textcolor{red}{不要記念}我幼年的罪愆和我的過犯，耶和華啊，求你因你的恩惠，按你的慈愛\textcolor{red}{記念}我。\\\hline
\end{tabular}
\end{table} 

\subsubsection{述说神对不同人的各种丰盛恩典 8-14}
耶和華是良善正直的，所以他必指示罪人走正路。
9 他必按公平引領謙卑人，將他的道教訓他們。
10 凡遵守他的約和他法度的人，耶和華都以慈愛、誠實待他。
11 耶和華啊，求你因你的名赦免我的罪，因為我的罪重大。
12 誰敬畏耶和華，耶和華必指示他當選擇的道路。
13 他必安然居住，他的後裔必承受地土。
14 耶和華與敬畏他的人親密，他必將自己的約指示他們。

\subsubsection{一个“必”七个“求”的祷告 15-22}
我的眼目時常仰望耶和華，\underline{因為}他\underline{必}將我的腳從網裡拉出來。
\underline{求}你轉向我，憐恤我，\underline{因為}我是孤獨困苦。我心裡的愁苦甚多，
\underline{求}你救我脫離我的禍患。
\underline{求}你看顧我的困苦我的艱難，赦免我一切的罪。
\underline{求}你察看我的仇敵，\underline{因為}他們人多，並且痛痛地恨我。
\underline{求}你保護我的性命，搭救我，使我不致羞愧，\underline{因為}我投靠你。
\underline{願}純全正直保守我，\underline{因為}我等候你。神啊，
\underline{求}你救贖以色列脫離他一切的愁苦！



\subsection{第26篇（大衛）： 自言行為純正，神前三求-申冤，察看，救贖 26:1-12}
\begin{table}[H]
\centering
\begin{tabular}{p{2.cm}|p{9.cm}|p{5.7500cm}}\hline\hline
\multicolumn{1}{c|}{求申冤1}&\multicolumn{1}{c}{求察看2-8}&\multicolumn{1}{c}{求救贖9-12}\\\hline
耶和華啊，求你為我申冤，因我向來行事純全，我又倚靠耶和華並不搖動。
&耶和華啊，求你察看我，試驗我，熬煉我的肺腑心腸。3 因為你的慈愛常在我眼前，我也按你的真理而行。4 我沒有和虛謊人同坐，也不與瞞哄人的同群。5 我恨惡惡人的會，必不與惡人同坐。6 耶和華啊，我要洗手表明無辜，才環繞你的祭壇，7 我好發稱謝的聲音，也要述說你一切奇妙的作為。8 耶和華啊，我喜愛你所住的殿和你顯榮耀的居所。
&不要把我的靈魂和罪人一同除掉，不要把我的性命和流人血的一同除掉，10 他們的手中有奸惡，右手滿有賄賂。至於我，卻要行事純全，求你救贖我，憐恤我。12 我的腳站在平坦地方，在眾會中我要稱頌耶和華。\\\hline
\end{tabular}
\end{table} 


\subsection{第27篇（大衛）： 以耶和華為拯救保障，心中無懼 27:1-14}
\subsubsection{神是拯救的宣告 1-3}
耶和華是我的亮光，是我的拯救，我還怕誰呢？耶和華是我性命的保障 a，我還懼誰呢？
2 那作惡的，就是我的仇敵，前來吃我肉的時候，就絆跌仆倒。
3 雖有軍兵安營攻擊我，我的心也不害怕；雖然興起刀兵攻擊我，我必仍舊安穩。
\subsubsection{曾向神求告并继续求告的一件事 4-6}
\begin{table}[H]
\centering
\begin{tabular}{p{5.75cm}|p{5.2500cm}|p{5.00cm}}\hline\hline
\multicolumn{1}{c|}{求告的一件事}&\multicolumn{1}{c|}{求告的原因}&\multicolumn{1}{c}{求告的结果}\\\hline
有一件事，我曾求耶和華，我仍要尋求，就是一生一世住在耶和華的殿中，瞻仰他的榮美，在他的殿裡求問。
&因為我遭遇患難，他必暗暗地保守我，在他亭子裡，把我藏在他帳幕的隱密處，將我高舉在磐石上。
&現在我得以昂首，高過四面的仇敵。我要在他的帳幕裡歡然獻祭，我要唱詩，歌頌耶和華。\\\hline
\end{tabular}
\end{table} 
\subsubsection{向神的祷告 7-12}
\begin{table}[H]
\centering
\begin{tabular}{p{2cm}|p{14.75cm}}\hline\hline
\multicolumn{1}{c|}{大衛}&\multicolumn{1}{c}{耶和華}\\\hline
&耶和華啊，我用聲音呼籲的時候，求你垂聽，並求你憐恤我，應允我。\\\hline
你說：「你們當尋求我的面。」&那時我心向你說：「耶和華啊，你的面我正要尋求。」不要向我掩面，不要發怒趕逐僕人，你向來是幫助我的。救我的神啊，不要丟掉我，也不要離棄我！
10 我父母離棄我，耶和華必收留我。
11 耶和華啊，求你將你的道指教我，因我仇敵的緣故引導我走平坦的路。
12 求你不要把我交給敵人，遂其所願，因為妄作見證的和口吐凶言的起來攻擊我。\\\hline
\end{tabular}
\end{table} 

\subsubsection{我在神面前的心愿 13-14}
我若不信在活人之地得見耶和華的恩惠，就早已喪膽了。
14 要等候耶和華！當壯膽，堅固你的心！我再說，要等候耶和華！


\subsection{第28篇（大衛）： 求神拯救脱离仇敌，并歌颂神垂听祷告 28:1-9}
\subsubsection{求耶和華垂聽其祈求 1-5}
\begin{table}[H]
\centering
\begin{tabular}{p{8.cm}|p{9.00cm}}\hline\hline
\multicolumn{1}{c|}{為自己祷告}&\multicolumn{1}{c}{為惡人祷告}\\\hline
耶和華啊，我要求告你！我的磐石啊，不要向我緘默！倘若你向我閉口，我就如將死的人一樣。2 我呼求你，向你至聖所舉手的時候，求你垂聽我懇求的聲音。不要把我和惡人並作孽的一同除掉，
&他們與鄰舍說和平話，心裡卻是奸惡。4 願你按著他們所做的，並他們所行的惡事待他們。願你照著他們手所做的待他們，將他們所應得的報應加給他們。5 他們既然不留心耶和華所行的和他手所做的，他就必毀壞他們，不建立他們。\\\hline
\end{tabular}
\end{table} 
\subsubsection{頌耶和華因允其求 6-9}
耶和華是應當稱頌的！因為他聽了我懇求的聲音。
7 耶和華是我的力量，是我的盾牌，我心裡倚靠他，就得幫助。所以我心中歡樂，我必用詩歌頌讚他。
8 耶和華是他百姓的力量，又是他受膏者得救的保障。
9 求你拯救你的百姓，賜福給你的產業，牧養他們，扶持他們，直到永遠。


\subsection{第29篇（大衛）： 歌頌耶和華荣耀威严和权柄 29:1-11}
\subsubsection{歌頌耶和華的荣耀 1-2}
神的眾子啊，你們要將榮耀、能力歸給耶和華，歸給耶和華！
2 要將耶和華的名所當得的榮耀歸給他，以聖潔的（為）裝飾敬拜耶和華！
\subsubsection{歌頌耶和華的聲音 3-9}
耶和華的聲音發在水上，榮耀的神打雷，耶和華打雷在大水之上。
4 耶和華的聲音大有能力，耶和華的聲音滿有威嚴。
5 耶和華的聲音震破香柏樹，耶和華震碎黎巴嫩的香柏樹。
6 他也使之跳躍如牛犢，使黎巴嫩和西連跳躍如野牛犢。
7 耶和華的聲音使火焰分岔。
8 耶和華的聲音震動曠野，耶和華震動加低斯的曠野。
9 耶和華的聲音驚動母鹿落胎，樹木也脫落淨光。凡在他殿中的，都稱說他的榮耀。
\subsubsection{歌頌耶和華的权柄音 10-11}
洪水氾濫之時，耶和華坐著為王；耶和華坐著為王，直到永遠。
11 耶和華必賜力量給他的百姓，耶和華必賜平安的福給他的百姓。


\subsection{第30篇（大衛）： 大衛獻殿時，讚美神的拯救恩典 30:1-12}
\subsubsection{赞美神的拯救 30:1-4}
大衛在獻殿的時候，作這詩歌。
1 耶和華啊，我要尊崇你，因為你曾提拔我，不叫仇敵向我誇耀。
2 耶和華我的神啊，我曾呼求你，你醫治了我。
3 耶和華啊，你曾把我的靈魂從陰間救上來，使我存活，不至於下坑。
4 耶和華的聖民哪，你們要歌頌他，稱讚他可記念的聖名。
\subsubsection{赞美神的恩典 30:5-12}
因為他的怒氣不過是轉眼之間，他的恩典乃是一生之久，一宿雖然有哭泣，早晨便必歡呼。6 至於我，我凡事平順，便說：「我永不動搖。」
7 耶和華啊，你曾施恩，叫我的江山穩固；你掩了面，我就驚惶。
8 耶和華啊，我曾求告你，我向耶和華懇求說：
9 「我被害流血，下到坑中，有什麼益處呢？塵土豈能稱讚你，傳說你的誠實嗎？
10 耶和華啊，求你應允我，憐恤我！耶和華啊，求你幫助我！」
11 你已將我的哀哭變為跳舞，將我的麻衣脫去，給我披上喜樂，
12 好叫我的靈（榮耀）歌頌你，並不住聲。耶和華我的神啊，我要稱謝你，直到永遠！


\subsection{第31篇（大衛）： 向神的九个“求“和赞美 31:1-24}
\subsubsection{向神的九个“求” 31:1-17}
耶和華啊，\textcolor{red}{求}你使我永不羞愧，憑你的公義搭救我 
\textcolor{red}{求}你側耳而聽，快快救我，做我堅固的磐石、拯救我的保障
因為你是我的巖石，我的山寨，所以\textcolor{red}{求}你為你名的緣故，引導我，指點我。
%\textcolor{red}{求}你救我脫離人為我暗設的網羅，因為你是我的保障。
\begin{table}[H]
\centering
\begin{tabular}{p{13.cm}|p{3.00cm}}\hline\hline
\multicolumn{2}{c}{\textcolor{red}{求}你救我脫離人為我暗設的網羅，因為你是我的保障。}\\\hline
%\multicolumn{1}{c|}{我}&\multicolumn{1}{c}{你}\\\hline
\textcolor{blue}{我}將我的靈魂交在你手裡，耶和華誠實的神啊，你救贖了我！\textcolor{blue}{我}恨惡那信奉虛無之神的人，我卻倚靠耶和華。\textcolor{blue}{我}要為你的慈愛高興歡喜，因為\textcolor{blue}{你}見過我的困苦，知道我心中的艱難。
&\textcolor{blue}{你}未曾把我交在仇敵手裡，\textcolor{blue}{你}使我的腳站在寬闊之處。\\\hline
\end{tabular}
\end{table} 
%耶和華啊，\textcolor{red}{求}你憐恤我，因為我在急難之中，
\begin{table}[H]
\centering
\begin{tabular}{p{5.5cm}|p{11.50cm}}\hline\hline
\multicolumn{2}{c}{耶和華啊，\textcolor{red}{求}你憐恤我，因為我在急難之中，}\\\hline
%\multicolumn{1}{c|}{身心也不安舒}&\multicolumn{1}{c}{四圍都是驚嚇}\\\hline
我的眼睛因憂愁而乾癟，連我的身心也不安舒。我的生命為愁苦所消耗，我的年歲為嘆息所曠廢，我的力量因我的罪孽衰敗，我的骨頭也枯乾。&我因一切敵人成了羞辱，在我的鄰舍跟前更甚。那認識我的都懼怕我，在外頭看見我的都躲避我。我被人忘記，如同死人，無人記念，我好像破碎的器皿。我聽見了許多人的讒謗，四圍都是驚嚇。他們一同商議攻擊我的時候，就圖謀要害我的性命。耶和華啊，我仍舊倚靠你，我說：「你是我的神！」\\\hline
\end{tabular}
\end{table} 

我終身的事在你手中，\textcolor{red}{求}你救我脫離仇敵的手和那些逼迫我的人。
\textcolor{red}{求}你使你的臉光照僕人，憑你的慈愛拯救我。
耶和華啊，\textcolor{red}{求}你叫我不致羞愧，因為我曾呼籲你。\textcolor{red}{求}你使惡人羞愧，使他們在陰間緘默無聲。


\subsubsection{向神的赞美 31:18-24}
\begin{table}[H]
\centering
\begin{tabular}{p{13.cm}|p{4.25cm}}\hline\hline
\multicolumn{1}{c|}{赞美神的作为 18-22}&\multicolumn{1}{c}{呼吁众人赞美23-24}\\\hline
那撒謊的人逞驕傲輕慢，出狂妄的話攻擊義人，願他的嘴啞而無言！
19 敬畏你、投靠你的人，你為他們所積存的，在世人面前所施行的恩惠是何等大呢！
20 你必把他們藏在你面前的隱密處，免得遇見人的計謀；你必暗暗地保守他們在亭子裡，免受口舌的爭鬧。
21 耶和華是應當稱頌的！因為他在堅固城裡，向我施展奇妙的慈愛。
22 至於我，我曾急促地說「我從你眼前被隔絕」，然而我呼求你的時候，你仍聽我懇求的聲音。
&耶和華的聖民哪，你們都要愛他！耶和華保護誠實人，足足報應行事驕傲的人。
24 凡仰望耶和華的人，你們都要壯膽，堅固你們的心！\\\hline
\end{tabular}
\end{table} 



\subsection{第32篇（大衛）： 罪得赦免乃為有福，馴服勸誡必有慈愛 32:1-11}
\begin{table}[H]
\centering
\begin{tabular}{p{11.cm}|p{6.cm}}\hline\hline
\multicolumn{1}{c|}{蒙主赦罪乃為有福1-7}&\multicolumn{1}{c}{馴服勸誡必有慈愛8-11}\\\hline
得赦免其過、遮蓋其罪的，這人是有福的!
2 凡心裡沒有詭詐，耶和華不算為有罪的，這人是有福的!
3 我閉口不認罪的時候，因終日唉哼而骨頭枯乾。
4 黑夜白日，你的手在我身上沉重，我的精液耗盡，如同夏天的乾旱。（細拉）
5 我向你陳明我的罪，不隱瞞我的惡。我說「我要向耶和華承認我的過犯」，你就赦免我的罪惡。（細拉）
6 為此，凡虔誠人都當趁你可尋找的時候禱告你，大水泛溢的時候，必不能到他那裡。
7 你是我藏身之處，你必保佑我脫離苦難，以得救的樂歌四面環繞我。（細拉）
&我要教導你，指示你當行的路，我要定睛在你身上勸誡你。
9 你不可像那無知的騾馬，必用嚼環轡頭勒住牠，不然就不能馴服。
10 惡人必多受苦楚，唯獨倚靠耶和華的，必有慈愛四面環繞他。
11 你們義人應當靠耶和華歡喜快樂，你們心裡正直的人都當歡呼!\\\hline
\end{tabular}
\end{table} 


\subsection{第33篇（大衛）： 頌讚耶和華之誠信，以耶和華為神之國有福 33:1-22}
\begin{table}[H]
\centering
\begin{tabular}{p{8.25cm}|p{9.cm}}\hline\hline
\multicolumn{1}{c|}{頌讚耶和華之誠信1-11}&\multicolumn{1}{c}{以耶和華為神之國有福12-22}\\\hline
義人哪，你們應當靠耶和華歡樂！正直人的讚美是合宜的。
2 你們應當彈琴稱謝耶和華，用十弦瑟歌頌他。
3 應當向他唱新歌，彈得巧妙，聲音洪亮。
4 因為耶和華的言語正直，凡他所做的盡都誠實。
5 他喜愛仁義、公平，遍地滿了耶和華的慈愛。
6 諸天藉耶和華的命而造，萬象藉他口中的氣而成。
7 他聚集海水如壘，收藏深洋在庫房。
8 願全地都敬畏耶和華，願世上的居民都懼怕他。
9 因為他說有，就有；命立，就立。
10 耶和華使列國的籌算歸於無有，使眾民的思念無有功效。
11 耶和華的籌算永遠立定，他心中的思念萬代常存。
&以耶和華為神的，那國是有福的！他所揀選為自己產業的，那民是有福的！
13 耶和華從天上觀看，他看見一切的世人，
14 從他的居所往外察看地上一切的居民。
15 他是那造成他們眾人心的，留意他們一切作為的。
16 君王不能因兵多得勝，勇士不能因力大得救。
17 靠馬得救是枉然的，馬也不能因力大救人。
18 耶和華的眼目看顧敬畏他的人和仰望他慈愛的人，
19 要救他們的命脫離死亡，並使他們在饑荒中存活。
20 我們的心向來等候耶和華，他是我們的幫助，我們的盾牌。
21 我們的心必靠他歡喜，因為我們向來倚靠他的聖名。
22 耶和華啊，求你照著我們所仰望你的，向我們施行慈愛。\\\hline
\end{tabular}
\end{table} 


\subsection{第34篇（大衛）： 赞美神的救恩，要离恶行善靠主有福 34:1-22}



\subsubsection{頌讚耶和華之救恩 34:1-11}
\begin{table}[H]
\centering
\begin{tabular}{p{5.25cm}|p{5.25cm}|p{6.5cm}}\hline\hline
\multicolumn{1}{c|}{你們和我當常常赞美神 1-3}&\multicolumn{1}{c|}{述说神的做为 4-7}&\multicolumn{1}{c}{呼吁众人一同投靠神 8-11}\\\hline
大衛在亞比米勒面前裝瘋，被他趕出去，就作這詩。
1 我要時時稱頌耶和華，讚美他的話必常在我口中。
2 我的心必因耶和華誇耀，謙卑人聽見，就要喜樂。
3 你們和我當稱耶和華為大，一同高舉他的名。
&我曾尋求耶和華，他就應允我，救我脫離了一切的恐懼。凡仰望他的便有光榮，他們的臉必不蒙羞。我這困苦人呼求，耶和華便垂聽，救我脫離一切患難。耶和華的使者在敬畏他的人四圍安營，搭救他們。
&你們要嘗嘗主恩的滋味，便知道他是美善。投靠他的人有福了!耶和華的\textcolor{red}{聖民}哪，你們當敬畏他，因敬畏他的一無所缺。少壯獅子還缺食忍餓，但尋求耶和華的，什麼好處都不缺。\textcolor{red}{眾弟子}啊，你們當來聽我的話，我要將敬畏耶和華的道教訓你們。\\\hline
\end{tabular}
\end{table} 


\subsubsection{离恶行善靠主有福 34:12-22}
\begin{table}[H]
\centering
\begin{tabular}{p{3.75cm}|p{13.25cm}}\hline\hline
\multicolumn{1}{c|}{离恶行善 12-14}&\multicolumn{1}{c}{義人的福气 15-22}\\\hline
有何人喜好存活，愛慕長壽，得享美福，就要禁止舌頭不出惡言，嘴唇不說詭詐的話。
14 要離惡行善，尋求和睦，一心追趕。
&耶和華的眼目看顧義人，他的耳朵聽他們的呼求。
16 耶和華向行惡的人變臉，要從世上除滅他們的名號。
17 義人呼求，耶和華聽見了，便救他們脫離一切患難。
18 耶和華靠近傷心的人，拯救靈性痛悔的人。
19 義人多有苦難，但耶和華救他脫離這一切，
20 又保全他一身的骨頭，連一根也不折斷。
21 惡必害死惡人，恨惡義人的必被定罪。
22 耶和華救贖他僕人的靈魂，凡投靠他的必不致定罪。\\\hline
\end{tabular}
\end{table} 


\subsection{第35篇（\underline{基督的顺服}大衛）：求耶和華拯救懲罰其敵 35:1-28}
耶和華啊，與我相爭的，求你與他們相爭；與我相戰的，求你與他們相戰。
2 拿著大小的盾牌，起來幫助我；
3 抽出槍來，擋住那追趕我的。求你對我的靈魂說：「我是拯救你的！」
4 願那尋索我命的，蒙羞受辱；願那謀害我的，退後羞愧。
5 願他們像風前的糠，有耶和華的使者趕逐他們。
6 願他們的道路又暗又滑，有耶和華的使者追趕他們。
7 因他們無故地為我暗設網羅，無故地挖坑要害我的性命。
8 願災禍忽然臨到他身上，願他暗設的網纏住自己，願他落在其中遭災禍。
我的心必靠耶和華快樂，靠他的救恩高興。
10 我的骨頭都要說：「耶和華啊，誰能像你，救護困苦人脫離那比他強壯的，救護困苦窮乏人脫離那搶奪他的？」

凶惡的見證人起來，盤問我所不知道的事。
12 他們向我以惡報善，使我的靈魂孤苦。
13 至於我，當他們有病的時候，我便穿麻衣，禁食刻苦己心，我所求的都歸到自己的懷中。
14 我這樣行，好像他是我的朋友，我的弟兄；我屈身悲哀，如同人為母親哀痛。
15 我在患難中，他們卻歡喜，大家聚集。我所不認識的那些下流人聚集攻擊我，他們不住地把我撕裂。
16 他們如同席上好戲笑的狂妄人，向我咬牙。


主啊，你看著不理要到幾時呢？求你救我的靈魂脫離他們的殘害，救我的生命(獨一者)脫離少壯獅子。
18 我在大會中要稱謝你，在眾民中要讚美你。
19 求你不容那無理與我為仇的向我誇耀，不容那無故恨我的向我擠眼。
20 因為他們不說和平話，倒想出詭詐的言語，害地上的安靜人。
21 他們大大張口攻擊我，說：「啊哈！啊哈！我們的眼已經看見了！」
22 耶和華啊，你已經看見了，求你不要閉口。主啊，求你不要遠離我。
23 我的神我的主啊，求你奮興醒起，判清我的事，申明我的冤。
24 耶和華我的神啊，求你按你的公義判斷我，不容他們向我誇耀。
25 不容他們心裡說：「阿哈！遂我們的心願了！」不容他們說：「我們已經把他吞了！」
26 願那喜歡我遭難的一同抱愧蒙羞，願那向我妄自尊大的披慚愧，蒙羞辱。
27 願那喜悅我冤屈得申(公義)的歡呼快樂，願他們常說：「當尊耶和華為大，耶和華喜悅他的僕人平安。」
28 我的舌頭要終日論說你的公義，時常讚美你。

\subsection{第36篇（\underline{基督的顺服}大衛）：惡人之罪惡,神之慈愛 36：1-12}
\subsubsection{惡人之罪惡}

耶和華的僕人大衛的詩，交於伶長。
1 惡人的罪過，在他心裡說：「我眼中不怕神。」
2 他自誇自媚，以為他的罪孽終不顯露，不被恨惡。
3 他口中的言語盡是罪孽詭詐，他與智慧善行已經斷絕。
4 他在床上圖謀罪孽，定意行不善的道，不憎惡惡事。

\subsubsection{神之慈愛}
5 耶和華啊，你的慈愛上及諸天，你的信實達到穹蒼。
6 你的公義好像高山，你的判斷如同深淵。耶和華啊，人民、牲畜你都救護。
7 神啊，你的慈愛何其寶貴，世人投靠在你翅膀的蔭下。
8 他們必因你殿裡的肥甘得以飽足，你也必叫他們喝你樂河的水。
9 因為在你那裡有生命的源頭，在你的光中我們必得見光。
10 願你常施慈愛給認識你的人，常以公義待心裡正直的人。
11 不容驕傲人的腳踐踏我，不容凶惡人的手趕逐我。
12 在那裡，作孽的人已經仆倒，他們被推倒，不能再起來。

\subsection{第37篇（大衛）： 不要因恶人而心懷不平而作惡，义人和恶人的结局 31:1-40}

\subsubsection{不要因作惡的而心懷不平，更不要因心懷不平而作惡 37:1-17}
\begin{table}[H]
\centering
\begin{tabular}{p{7cm}|p{10.00cm}}\hline\hline
\multicolumn{1}{c|}{不要因作惡的而心懷不平 1-7}&\multicolumn{1}{c}{不要因心懷不平而作惡 8-17}\\\hline
\textcolor{red}{不要}為作惡的心懷不平，也\textcolor{red}{不要}向那行不義的生出嫉妒。
2 因為他們如草快被割下，又如青菜快要枯乾。
3 你\textcolor{red}{當}倚靠耶和華而行善，住在地上，以他的信實為糧。
4 又\textcolor{red}{要}以耶和華為樂，他就將你心裡所求的賜給你。
5 \textcolor{red}{當}將你的事交託耶和華，並倚靠他，他就必成全。
6 他\textcolor{red}{要}使你的公義如光發出，使你的公平明如正午。
7 你\textcolor{red}{當}默然倚靠耶和華，耐性等候他，\textcolor{red}{不要}因那道路通達的和那惡謀成就的心懷不平。
&當止住怒氣，離棄憤怒，不要心懷不平，以致作\textcolor{red}{惡}。
9 因為作\textcolor{red}{惡}的必被剪除，唯有\underline{等候耶和華的必承受地土}。
10 還有片時，\textcolor{red}{惡人}要歸於無有，你就是細察他的住處，也要歸於無有。
11 但\underline{謙卑人必承受地土，以豐盛的平安為樂}。
12 \textcolor{red}{惡人}設謀害\textcolor{red}{義人}，又向他咬牙。
13 主要笑他，因見他受罰的日子將要來到。
14 \textcolor{red}{惡人}已經弓上弦，刀出鞘，要打倒困苦窮乏的人，要殺害行動正直的人。
15 他們的刀必刺入自己的心，他們的弓必被折斷。
16 一個\textcolor{red}{義人}所有的雖少，\underline{強過許多\textcolor{red}{惡人}的富餘}。
17 因為\textcolor{red}{惡人}的膀臂必被折斷，但\underline{耶和華是扶持\textcolor{red}{義人}}。\\\hline
\end{tabular}
\end{table} 

\begin{itemize}
\itemsep-.45em 
\item 不要因作惡的而心懷不平 1-7
\begin{itemize}
\itemsep-.45em 
\item 面对恶人恶事的错误做法（四个因）及恶人的结局
\begin{itemize}
\itemsep-.45em 
\item \underline{因}作惡的心懷不平，如草快被割下，又如青菜快要枯乾
\item \underline{因}那行不義的生出嫉妒	
\item \underline{因}那道路通達的心懷不平	
\item \underline{因}那那惡謀成就的心懷不平	
\end{itemize}
\item 面对恶人恶事的正确做法（五个當）及结局：
\begin{itemize}
\itemsep-.45em 
\item \underline{當}倚靠耶和華而行善	
\item \underline{（當）}住在地上以他的信實為糧	
\item \underline{（當）}以耶和華為樂，他就將你心裡所求的賜給你
\item \underline{（當）}將你的事交託耶和華，並倚靠他，他就必成全，使你的公義如光發出，使你的公平明如正午
\item \underline{當}默然倚靠耶和華，耐性等候他	
\end{itemize}
\end{itemize}
\item 不要因心懷不平而作惡 8-17
\end{itemize}

\subsubsection{義人和惡人的两个见证 37:18-40}
\begin{table}[H]
\centering
\begin{tabular}{p{7.5cm}|p{9.5cm}}\hline\hline
\multicolumn{1}{c|}{義人终将承受地土 18-27}&\multicolumn{1}{c}{惡人住处和后裔被剪除 28-40}\\\hline
耶和華知道完全人的日子，他們的產業要存到永遠。
他們在急難的時候不致羞愧，在饑荒的日子必得飽足。
惡人卻要滅亡，耶和華的仇敵要像羊羔的脂油(像草地的華美)，他們要消滅，要如煙消滅。
惡人借貸而不償還，義人卻恩待人，並且施捨。
蒙耶和華賜福的必承受地土，被他咒詛的必被剪除。
義人的腳步被耶和華立定，他的道路耶和華也喜愛。
他雖失腳，也不至全身仆倒，因為耶和華用手攙扶他(攙扶他的手)。

\textcolor{blue}{我從前年幼，現在年老，卻未見過義人被棄，也未見過他的後裔討飯。
他終日恩待人，借給人，他的後裔也蒙福}。
你當離惡行善，就可永遠安居。
&因為耶和華喜愛公平，不撇棄他的聖民，他們永蒙保佑，但惡人的後裔必被剪除。
義人必承受地土，永居其上。
義人的口談論智慧，他的舌頭講說公平。
神的律法在他心裡，他的腳總不滑跌。
惡人窺探義人，想要殺他。
耶和華必不撇他在惡人手中，當審判的時候，也不定他的罪。
你當等候耶和華，遵守他的道，他就抬舉你，使你承受地土，惡人被剪除的時候，你必看見。

\textcolor{blue}{我見過惡人大有勢力，好像一根青翠樹在本土生發。
有人從那裡經過，不料，他沒有了；我也尋找他，卻尋不著。
你要細察那完全人，觀看那正直人，因為和平人有好結局。
至於犯法的人，必一同滅絕，惡人終必剪除}。
但義人得救，是由於耶和華，他在患難時做他們的營寨。
耶和華幫助他們，解救他們。他解救他們脫離惡人，把他們救出來，因為他們投靠他。\\\hline
\end{tabular}
\end{table} 


\subsubsection{教導}
\begin{itemize}
\itemsep-.45em 
\item 这首诗歌的作者讲出了因为遭遇恶人或者恶事而采取的两种错误做法。
人行恶的原因有多种，比如：（1）有的人行恶是因为眼中没有神（诗54：3）；（2）有的人行恶是知道有神，但因为心中轻慢神而行恶（诗10：13）；（3）这首诗篇里谈到另外一种行恶的原因：就是有的人看到恶人计谋屡屡得逞，甚至兴盛发达，因而心中不平，这本身就已经偏离了神的道了，但这种人不仅仅停留在心中的恼恨和不平，而是因着心中不平，产生愤怒，也开始行恶。这首诗歌详细分析了第三种人犯罪的步骤，就是：

\begin{itemize}
\itemsep-.45em 
\item 為作惡的心懷不平,心中生出嫉妒
作者在诗歌中首先给我们分享了我们所嫉妒的恶人的结局是：生命如青菜般的短暂，结局是被割下，也就是灭亡。提醒我们不要因恶人今天的得意而羡慕嫉妒他们。作者又从正面给出了当心中涌出不平和嫉妒时，正确处理我们错误情绪的五种做法：就是（1）住在地上要依靠神而行善。“依靠神”的意思就是要按照神的标准来行善事，而不是用自己眼中的善恶标准。当我们照着神的是非标准行事为人时，神的话必保守我们象一颗树，不枯干，按时候结果子，凡我们所做的，尽都顺利（诗1：3）。（2）我们要每天要依靠神的信实而活，当我们依靠神的信实时，就不会羡慕或者嫉妒别人拥有，而我们目前暂时还没有，或者我们一生都不需要的东西。（3）还要以神为乐，如果我们以神为乐，地上的世人所有的虚荣自然不会再对我们有吸引力，并且作者告诉我们，我们所需要的神必赐给我们。（4）将我们所要的交托依靠神，诗人指出，神要成全，并使我们的“公義如光發出，公平明如正午”。神在这里的应许不但成就我们所交托的，而且还将我们因着信靠他，而从神得到的义，也要显明出来！（5）我们當默然倚靠，耐性等候神。有时候我们交托给神，并不一定我们马上就能得到，但不代表神不成就，所以我们要学习化时间等候神来成就。
\item 因为心懷不平，以致作惡。
作者在这里首先告诉我们避免因为心懷不平以致作惡的第一步就是“當止住怒氣，離棄憤怒”。这样做就可以从。作者接下来又用笔墨给读者分析了，因为心懷不平以致作惡的恶人的结局，和不去行恶的义人的结局。透过这两给对比，提醒我们不要效法恶人。
\end{itemize}

\item 大卫写这首诗歌的时候应该已经是老年了，他用自己的一生与神同行的经历为神做了关于义人的见证：那就是大卫一生“未見過義人被棄，也未見過他的後裔討飯”（37：25-34）。義人的後裔也蒙福。
\item 从本诗可以看出大卫是一个察验神作为的人。他不但查验義人，而且也查验恶人的结局，恶人虽然如一颗大的青翠樹在本土生旺，但终必被剪除（37：35-40）。
\end{itemize}

\subsection{第38篇（大衛）： 哀陳苦況求主矜憐 38:1-22}
\subsubsection{因罪過被神責備身患重病 1-10}
耶和華啊，求你不要在怒中責備我，不要在烈怒中懲罰我！
2 因為你的箭射入我身，你的手壓住我。
3 因你的惱怒，我的肉無一完全；因我的罪過，我的骨頭也不安寧。
4 我的罪孽高過我的頭，如同重擔叫我擔當不起。
5 因我的愚昧，我的傷發臭流膿。
6 我疼痛，大大蜷曲，終日哀痛。
7 我滿腰是火，我的肉無一完全。
8 我被壓傷，身體疲倦；因心裡不安，我就唉哼。
9 主啊，我的心願都在你面前，我的嘆息不向你隱瞞。
10 我心跳動，我力衰微，連我眼中的光也沒有了。
\subsubsection{親友遠离仇敵作對，我要認罪因罪憂愁，求神幫助 11-22}
\begin{table}[H]
\centering
\begin{tabular}{p{7cm}|p{10.25cm}}\hline\hline
\multicolumn{1}{c|}{良朋密友、親戚本家、仇敵}&\multicolumn{1}{c}{大衛}\\\hline
我的良朋密友因我的災病都躲在旁邊站著，

我的親戚本家也遠遠地站立。

那尋索我命的設下網羅，

那想要害我的口出惡言，終日思想詭計。
&但我如聾子不聽，像啞巴不開口。我如不聽見的人，口中沒有回話。耶和華啊，我仰望你！主我的神啊，你必應允我！我曾說：「恐怕他們向我誇耀，我失腳的時候，他們向我誇大。」我幾乎跌倒，我的痛苦常在我面前。我要承認我的罪孽，我要因我的罪憂愁。\\\hline
但我的仇敵又活潑又強壯，無理恨我的增多了。以惡報善的與我作對，因我是追求良善。&耶和華啊，求你不要撇棄我！我的神啊，求你不要遠離我！拯救我的主啊，求你快快幫助我！\\\hline
\end{tabular}
\end{table} 


\subsection{第39篇（大衛）：世事虛幻人生苦短，求神寬容唯主是望 39:1-13}
\subsubsection{惡人面前謹言慎行 1-2}
大衛的詩，交於伶長耶杜頓。
1 我曾說：「我要謹慎我的言行，免得我舌頭犯罪。惡人在我面前的時候，我要用嚼環勒住我的口。」
2 我默然無聲，連好話也不出口，我的愁苦就發動了。
\subsubsection{心中愁苦如火燒，便用口向神祷告 3-13}
3 我的心在我裡面發熱，我默想的時候，火就燒起，我便用舌頭說話。
4 耶和華啊，求你叫我曉得我身之終，我的壽數幾何，叫我知道我的生命不長。
5 你使我的年日窄如手掌，我一生的年數，在你面前如同無有。各人最穩妥的時候，真是全然虛幻！（細拉）
6 世人行動實係幻影，他們忙亂真是枉然，積蓄財寶不知將來有誰收取。
7 主啊，如今我等什麼呢？我的指望在乎你。
8 求你救我脫離一切的過犯，不要使我受愚頑人的羞辱。
9 因我所遭遇的是出於你，我就默然不語。
10 求你把你的責罰從我身上免去，因你手的責打，我便消滅。
11 你因人的罪惡，懲罰他的時候，叫他的笑容(所喜愛的)消滅，如衣被蟲所咬。世人真是虛幻！（細拉）
12 耶和華啊，求你聽我的禱告，留心聽我的呼求，我流淚，求你不要靜默無聲。因為我在你面前是客旅，是寄居的，像我列祖一般。
13 求你寬容我，使我在去而不返之先，可以力量復原。

\subsection{第40篇（\underline{基督的顺服}大衛）： 蒙主救拔脫離苦難???? 40:1-17}
\subsubsection{神所行的奇事 1-5}
大衛的詩，交於伶長。
1 我曾耐性等候耶和華，他垂聽我的呼求。
2 他從禍坑裡，從淤泥中，把我拉上來，使我的腳立在磐石上，使我腳步穩當。
3 他使我口唱新歌，就是讚美我們神的話。許多人必看見而懼怕，並要倚靠耶和華。
4 那倚靠耶和華，不理會狂傲和偏向虛假之輩的，這人便為有福！
5 耶和華我的神啊，你所行的奇事，並你向我們所懷的意念甚多，不能向你陳明。若要陳明，其事不可勝數。
\subsubsection{预言基督以身体做贖罪祭 6-10}
6 祭物和禮物你不喜悅，你已經開通我的耳朵，燔祭和贖罪祭非你所要。
7 那時我說：「看哪，我來了，我的事在經卷上已經記載了。
8 我的神啊，我樂意照你的旨意行，你的律法在我心裡。」
9 我在大會中宣傳公義的佳音，我必不止住我的嘴唇，耶和華啊，這是你所知道的。
10 我未曾把你的公義藏在心裡，我已陳明你的信實和你的救恩，我在大會中未曾隱瞞你的慈愛和誠實。
\subsubsection{向神的祷告和心愿 11-17}
11 耶和華啊，求你不要向我止住你的慈悲，願你的慈愛和誠實常常保佑我。
12 因有無數的禍患圍困我，我的罪孽追上了我，使我不能昂首。這罪孽比我的頭髮還多，我就心寒膽戰。
13 耶和華啊，求你開恩搭救我！耶和華啊，求你速速幫助我！
14 願那些尋找我要滅我命的一同抱愧蒙羞，願那些喜悅我受害的退後受辱。
15 願那些對我說「啊哈！啊哈！」的，因羞愧而敗亡。
16 願一切尋求你的因你高興歡喜，願那些喜愛你救恩的常說：「當尊耶和華為大！」
17 但我是困苦窮乏的，主仍顧念我。你是幫助我的，搭救我的，神啊，求你不要耽延！

\subsection{第41篇（\underline{基督被出卖}大衛）：病重時仇敵朋友惡言議論，求神憐恤  41:1-13}
\begin{table}[H]
\centering
\begin{tabular}{p{5.75cm}|p{6.00cm}|p{5.00cm}}\hline\hline
\multicolumn{1}{c|}{神的扶持1-4}&\multicolumn{1}{c|}{仇敵朋友用惡言議論5-9}&\multicolumn{1}{c}{求神憐恤10-13}\\\hline
1 眷顧貧窮的有福了！他遭難的日子，耶和華必搭救他。
2 耶和華必保全他，使他存活，他必在地上享福——求你不要把他交給仇敵，遂其所願。
3 他病重在榻，耶和華必扶持他；他在病中，你必給他鋪床。
4 我曾說：「耶和華啊，求你憐恤我，醫治我，因為我得罪了你。」
&5 我的仇敵用惡言議論我說：「他幾時死，他的名才滅亡呢？」
6 他來看我，就說假話；他心存奸惡，走到外邊才說出來。
7 一切恨我的，都交頭接耳地議論我，他們設計要害我。
8 他們說：「有怪病貼在他身上，他已躺臥，必不能再起來。」
9 連我知己的朋友，我所倚靠吃過我飯的，也用腳踢我。
&10 耶和華啊，求你憐恤我，使我起來，好報復他們。
11 因我的仇敵不得向我誇勝，我從此便知道你喜愛我。
12 你因我純正，就扶持我，使我永遠站在你的面前。
13 耶和華以色列的神是應當稱頌的，從亙古直到永遠！阿們，阿們。\\\hline
\end{tabular}
\end{table} 

















